\chapter{操作配方}

本章示例面向普通项目根目录。不要把 CCB 源码 checkout \src{/home/bfly/yunwei/ccb_source} 当作运行验证目录；源码验证应使用专门外部测试项目。

\section{创建最小项目配置}

\begin{lstlisting}
cmd | archi:codex
\end{lstlisting}

验证：

\begin{lstlisting}
ccb config validate
\end{lstlisting}

启动：

\begin{lstlisting}
ccb
\end{lstlisting}

\section{两 agent 协作}

\begin{lstlisting}
cmd | archi:codex | reviewer:claude
\end{lstlisting}

提交任务：

\begin{lstlisting}
ccb ask archi analyze current architecture
ccb ask reviewer review archi's result
\end{lstlisting}

\section{使用 windows topology}

\begin{lstlisting}
version = 2
entry_window = "main"

[windows]
main = "archi:codex | reviewer:claude"
ops = "ops:codex"

[ui.sidebar]
mode = "every_window"
width = "15%"
\end{lstlisting}

\section{添加 rich workbench}

\begin{lstlisting}
version = 2
entry_window = "main"

[windows]
main = "archi:codex, rich"
\end{lstlisting}

安装或更新 rich：

\begin{lstlisting}
ccb update rich
\end{lstlisting}

\section{安装并使用 Archi role}

\begin{lstlisting}
ccb roles install agentroles.archi
ccb roles add agentroles.archi:codex
ccb config validate
ccb reload --dry-run
ccb reload
\end{lstlisting}

\section{长报告工作流}

\begin{lstlisting}
ccb ask --artifact-reply archi write detailed architecture report
ccb trace <job_id>
\end{lstlisting}

\section{父任务请求子任务证据}

在 active CCB task 中：

\begin{lstlisting}
ccb ask --chain --artifact-reply archi collect source evidence
\end{lstlisting}

不要用普通 nested ask 等待结果。

\section{诊断失败任务}

\begin{lstlisting}
ccb trace job_xxx
ccb watch job_xxx
ccb queue --detail archi
ccb inbox --detail talk1
ccb doctor storage
\end{lstlisting}

\section{重试失败 attempt}

\begin{lstlisting}
ccb trace job_xxx
ccb retry job_xxx
ccb watch <new_job_id>
\end{lstlisting}

\section{重新提交旧 message}

\begin{lstlisting}
ccb trace msg_xxx
ccb resubmit msg_xxx
\end{lstlisting}

\section{清理已读 reply}

\begin{lstlisting}
ccb inbox talk1
ccb ack talk1
\end{lstlisting}

如果 ack 被拒绝，按提示使用当前 head inbound event id。
